\subsection{Course-of-values induction}
The course-of-values induction may be rephrased as follows:
\begin{thm}[Course-of-values induction]
Let $P(n)$ be a property about natural number $n$. Let $Q(m)$ be a property
about natural number $m$ such that it holds iff $P(n)$ holds for all $n \le m$.

If $P(0)$ and $P(m+1)$ follows from $Q(m)$, then $P(n)$ is true for all $n$. 
\end{thm}
\begin{proof}
The proof is quite simple.

First, prove $Q(0)$ (which is the same is $P(0)$), and then see the
obvious fact that $P(m+1)$ with $Q(m)$ deduces $Q(m+1)$. Then, we apply the
induction principle on $Q$, to conclude $Q(m)$ for all $m$. Clearly, $Q(m)$
implies $P(m)$, and we are done.
\end{proof}

It may also be rephrased more simply, where $Q(m)$ holds whenever $P(n)$ holds
for all $n < m$. Then, $Q(0)$ is vacuously true, and the rest is the same.
